\documentclass[conference]{IEEEtran}
\IEEEoverridecommandlockouts

\usepackage{cite}
\usepackage{amsmath,amssymb,amsfonts}
\usepackage{algorithmic}
\usepackage{graphicx}
\usepackage{textcomp}
\usepackage{xcolor}
\usepackage{multirow}
\usepackage{booktabs}
\def\BibTeX{{\rm B\kern-.05em{\sc i\kern-.025em b}\kern-.08em
    T\kern-.1667em\lower.7ex\hbox{E}\kern-.125emX}}

\begin{document}

\title{Joint Routing and Wavelength Assignment with Deep Reinforcement Learning for Quantum-Classical Coexistence in WDM Optical Networks}

\author{\IEEEauthorblockN{Yuting Wang}
\IEEEauthorblockA{\textit{School of Information and Communication Engineering} \\
\textit{Beijing University of Posts and Telecommunications}\\
Beijing, China \\
wangyuting007@bupt.edu.cn}
}

\maketitle

\begin{abstract}
Quantum key distribution (QKD) over existing wavelength-division multiplexing (WDM) optical networks requires solving a joint routing and wavelength assignment (RWA) problem under dual capacity constraints: classical data bandwidth and quantum key generation rate. The extreme power disparity between classical (milliwatt-level) and quantum (single-photon-level) signals means that nonlinear fiber noise---four-wave mixing (FWM) and spontaneous Raman scattering (SpRS)---from classical channels can severely degrade QKD performance, making the key capacity of each link dependent on the wavelength-level allocation of classical traffic. This paper formulates the QKD-aware RWA problem in WDM networks and proposes a two-layer solution framework: at the routing layer, we evaluate four path-selection strategies with increasing degrees of QKD awareness (min-hop, min-distance, key-capacity-aware, and dual-capacity-aware); at the wavelength assignment layer, we introduce a deep Q-network (DQN) agent that learns to select wavelengths to minimize FWM and SpRS crosstalk onto quantum channels. The DQN agent is trained via Stable-Baselines3 on a single-core WDM environment that integrates discrete single-frequency FWM and SpRS noise models with a dynamic Poisson traffic simulator. Simulation results on the NSFNET topology show that dual-capacity-aware routing reduces QKD blocking by up to 35\% compared to min-hop routing, and the DQN wavelength selection outperforms first-fit by up to 28\% in aggregate secret key rate under tight channel spacing (50~GHz). These results establish an integrated routing-and-allocation framework for reliable QKD coexistence in dense WDM networks.
\end{abstract}

\begin{IEEEkeywords}
Quantum key distribution, routing and wavelength assignment, deep reinforcement learning, four-wave mixing, Raman scattering, WDM optical networks
\end{IEEEkeywords}

\section{Introduction}

Quantum key distribution (QKD) \cite{b1} promises information-theoretically secure symmetric keys, with security rooted in quantum mechanics rather than computational hardness assumptions. As QKD matures toward practical wide-area deployment, integrating quantum channels into existing wavelength-division multiplexing (WDM) optical networks has emerged as the most economically viable path---it avoids the prohibitive cost of dedicated dark fibers by co-propagating quantum and classical signals in the same fiber \cite{b2}.

This coexistence, however, introduces a fundamental physical-layer challenge that propagates upward to the network layer. Classical WDM channels launch at 0--10~dBm per wavelength, while quantum signals carry less than one photon per pulse on average ($\sim -100$~dBm). The resulting $>$80~dB power gap makes QKD receivers exquisitely sensitive to nonlinear crosstalk from classical channels. Once linear crosstalk mechanisms are suppressed below $-110$~dB, two fiber nonlinear effects dominate \cite{b3}: \emph{four-wave mixing} (FWM), a Kerr-effect-driven process wherein classical channel triplets generate new frequency components that can land directly in quantum receiver passbands, and \emph{spontaneous Raman scattering} (SpRS), broadband inelastic photon scattering that transfers classical pump power across the entire C-band.

Critically, both FWM and SpRS noise depend on \emph{which wavelengths} carry classical traffic. This couples the physical-layer impairment to the network-layer routing and wavelength assignment (RWA) decisions: a routing decision that appears feasible under nominal key capacity may become infeasible if wavelength assignment places classical channels in FWM-resonant positions relative to quantum channels. The resulting network resource allocation problem is \textbf{dual-constrained}: each fiber link provides classical data capacity (Gbps) and QKD key capacity (kbps), where the latter depends nonlinearly on the classical wavelength occupancy.

From the perspective of communication network theory, this problem instantiates a \emph{reliability-constrained resource allocation} task. A lightpath is considered reliable if it maintains positive secret key rate (SKR) for every quantum channel under the FWM and SpRS noise generated by all co-propagating classical channels. The reliability constraint creates a structured coupling between the routing layer (which paths to use) and the wavelength assignment layer (which frequency slots to occupy on each link), making the joint optimization NP-hard.

Prior work has addressed routing and wavelength assignment for QKD networks separately. Zhang et al.\ \cite{b4} formulated an integer linear program for joint routing, channel, key-rate, and time-slot assignment but assumed fixed per-wavelength key capacities. Sun et al.\ \cite{b5} and Bahrani et al.\ \cite{b6} proposed wavelength assignment heuristics (equal-spacing and Raman-optimized, respectively) for static channel plans but did not integrate dynamic traffic or routing. On the machine learning front, deep reinforcement learning (DRL) has shown promise for classical RWA \cite{b7}, but its application to QKD-aware resource allocation remains largely unexplored.

In this paper, we make the following contributions:
\begin{enumerate}
    \item We formulate the QKD-aware RWA problem for single-fiber WDM coexistence networks, with link key capacities that depend on per-wavelength classical occupancy through discrete single-frequency FWM and SpRS noise models.
    \item We design a two-layer solution framework: at the routing layer, four strategies with increasing QKD awareness (min-hop, min-distance, key-capacity-aware, dual-capacity-aware) select end-to-end paths; at the wavelength assignment layer, both a first-fit heuristic and a DQN-based learned policy allocate wavelength slots.
    \item We train the DQN agent using Stable-Baselines3 on a custom Gymnasium environment that integrates the discrete FWM/SpRS noise models with Poisson traffic generation, and demonstrate that the learned policy outperforms first-fit wavelength assignment.
\end{enumerate}

\section{Network Modeling}

\subsection{WDM Network Topology and Traffic Model}

We model the optical network as a directed graph $G = (V, E)$, where vertices represent optical switching nodes and edges represent fiber spans. We use the NSFNET 14-node, 21-edge topology as the reference network, with edge lengths ranging from 52.5~km to 300~km. Each edge carries $W$ WDM wavelengths on the ITU-T G.694.1 C-band grid, with center frequencies
\begin{equation}
f(w) = f_{\text{start}} + (w - 1) \cdot \Delta f_{\text{ch}}, \quad w \in \{1, \ldots, W\},
\label{eq:itu_grid}
\end{equation}
where $f_{\text{start}} = 190.1$~THz, $\Delta f_{\text{ch}} = 100$~GHz for the standard grid (flexible down to 25~GHz), and $W = 61$ covers the full C-band.

Service requests arrive following a Poisson process with rate $\lambda$ and hold for exponentially distributed durations with mean $1/\mu$, yielding offered load $\rho = \lambda / \mu$ Erlangs. Each request $r$ specifies a source-destination pair $(s_r, d_r)$, a classical bandwidth demand $b_r$ (Gbps), and a key consumption demand $k_r$ (kbps).

\subsection{Dual-Capacity Resource Model}

Each edge $e \in E$ provides two types of capacity:
\begin{itemize}
    \item \textbf{Classical capacity} $C_e^{\text{cl}}$ (Gbps): determined by the modulation format and spectral efficiency of active classical wavelengths.
    \item \textbf{Key capacity} $K_e(\mathbf{x}_e)$ (kbps): the aggregate secret key rate supported by quantum channels on edge $e$, which depends on the binary occupancy vector $\mathbf{x}_e \in \{0,1\}^W$ indicating which wavelengths carry classical traffic on that edge.
\end{itemize}

The key capacity is computed from the SKR model after evaluating the FWM and SpRS noise generated by all active classical wavelengths. When the classical occupancy of a link changes (due to an arrival or departure), the key capacity is recalculated.

\subsection{Discrete Single-Frequency Noise Models}

We adopt the discrete single-frequency approach, treating each classical channel as a monochromatic pump at its ITU grid center frequency with launch power $P_{\text{ch}}$.

\subsubsection{FWM Noise}

For a quantum channel at frequency $f_q$, the forward FWM noise power at fiber length $L$ is computed by enumerating all classical channel triplets $(i, j, k)$ satisfying the frequency-matching condition $f_i + f_j - f_k = f_q$ \cite{b8}:
\begin{equation}
P_{\text{FWM}}^{\text{fwd}}(f_q, L) = \frac{\gamma^2 e^{-\alpha_q L}}{9} \sum_{\substack{(i,j,k) \\ f_i+f_j-f_k=f_q}} D_{ijk}^2 \cdot \eta_{ijk}(L) \cdot P_i P_j P_k,
\label{eq:fwm}
\end{equation}
where $\gamma$ is the fiber nonlinear coefficient, $D_{ijk} \in \{3, 6\}$ is the degeneracy factor, and $\eta_{ijk}(L)$ is the FWM efficiency:
\begin{equation}
\eta_{ijk}(L) = \left|\frac{1 - \exp[-\alpha L + j\Delta\beta_{ijk} L]}{\alpha - j\Delta\beta_{ijk}}\right|^2,
\label{eq:eta}
\end{equation}
with phase mismatch expanded to second-order dispersion:
\begin{equation}
\Delta\beta_{ijk} \approx \frac{2\pi\lambda^2}{c} (f_i - f_k)(f_j - f_k) \left[D_c + \frac{\lambda^2}{2c} \frac{dD_c}{d\lambda}(f_i + f_j - 2f_0)\right].
\label{eq:delta_beta}
\end{equation}

Backward FWM arises from Rayleigh backscattering of the forward-generated FWM power integrated along the fiber span. The total FWM noise power is $P_{\text{FWM}} = P_{\text{FWM}}^{\text{fwd}} + P_{\text{FWM}}^{\text{bwd}}$.

\subsubsection{SpRS Noise}

The SpRS noise from a classical pump at $f_p$ to a quantum channel at $f_q$ is computed via the Raman cross-section with phonon population correction \cite{b9}:
\begin{equation}
P_{\text{SpRS}} = \sum_{p \in \mathcal{C}} P_p \cdot \sigma(f_p, f_q) \cdot [n_{\text{th}}(\Delta f) + \Theta(\Delta f)] \cdot \mathcal{P}(f_p, f_q, L),
\label{eq:sprs}
\end{equation}
where $\Delta f = f_p - f_q$, $n_{\text{th}}(\Delta f) = [\exp(h|\Delta f|/kT) - 1]^{-1}$ is the phonon occupation factor, $\Theta(\Delta f) = 1$ for Stokes ($\Delta f > 0$) and $0$ otherwise, $\sigma(f_p, f_q)$ incorporates the measured 92-point GNPy Raman gain coefficient $g_R(|\Delta f|)$, and $\mathcal{P}$ is the distributed propagation factor distinguishing forward and backward scattering.

\subsection{Secret Key Rate Computation}

The noise photon probability per detector gate is
\begin{equation}
p_{\text{noise}} = \frac{P_{\text{FWM}} + P_{\text{SpRS}}}{2 h f_q} \cdot \tau_{\text{SPD}} \cdot \eta_{\text{SPD}} \cdot \text{IL}_{\text{post}}.
\label{eq:p_noise}
\end{equation}

The SKR under the three-intensity decoy-state BB84 protocol with finite-key corrections ($N = 10^{10}$ pulses) is \cite{b10}:
\begin{equation}
R_{\text{SKR}} = q \cdot \left\{ -Q_\mu f_{\text{EC}} H_2(E_\mu) + p_1^\mu Y_1^L [1 - H_2(e_1^U)] \right\} - \Delta_{\text{finite}},
\label{eq:skr}
\end{equation}
with protocol parameters $\mu = 0.75$ (signal), $\nu = 0.2$ (decoy), $q = 1/2$, $f_{\text{EC}} = 1.16$, detector efficiency $\eta_{\text{SPD}} = 0.25$, dark count probability $p_{\text{dark}} = 10^{-6}$, and gate width $\tau_{\text{SPD}} = 1$~ns. The insertion loss $\text{IL} = 6$~dB is applied after the shared fiber span.

A link $e$ is \emph{reliable} if $R_{\text{SKR}} > 0$ for all quantum channels on that link given the current classical occupancy $\mathbf{x}_e$. The total key capacity of link $e$ is $K_e(\mathbf{x}_e) = N_q \cdot R_{\text{SKR}} \cdot R_{\text{rep}}$, where $N_q$ is the number of parallel quantum channels and $R_{\text{rep}} = 50$~MHz is the pulse repetition rate.

Figure~\ref{fig:skr_distance} shows the finite-key SKR vs.\ fiber distance for the default protocol parameters, computed using the self-contained \texttt{qkd\_routing/physics.py} module. At zero noise ($p_{\text{noise}} = 0$), the SKR remains positive beyond 300~km. With realistic SpRS noise from adjacent classical channels, the achievable distance shortens significantly, motivating the joint routing and wavelength assignment optimization.

\begin{figure}[htbp]
\centerline{\includegraphics[width=\columnwidth]{results/skr_vs_distance.png}}
\caption{Finite-key decoy-state BB84 SKR vs.\ fiber distance for 0–8 co-propagating classical channels at $-25$~dBm/channel. Raman scattering from each additional classical channel degrades the achievable distance, motivating QKD-aware resource allocation.}
\label{fig:skr_distance}
\end{figure}

\section{Algorithm Design}

\subsection{Problem Formulation}

The joint routing and wavelength assignment (RWA) problem for QKD coexistence is:
\begin{equation}
\begin{aligned}
\underset{\text{path } p, \text{ wavelength } w}{\text{maximize}} \quad & U(p, w; \mathbf{X}) \\
\text{subject to} \quad & w \text{ is free on all edges of } p, \\
& K_e(\mathbf{x}_e + \mathbf{1}_w) > 0, \quad \forall e \in p,
\end{aligned}
\label{eq:rwa}
\end{equation}
where $\mathbf{X}$ is the global wavelength occupancy state, $\mathbf{1}_w$ is the one-hot vector for wavelength $w$, and $U(\cdot)$ is a network utility function. The second constraint ensures that adding a classical channel at wavelength $w$ does not kill key generation on any link along the path.

We decompose the problem into two stages: (1) \emph{routing} selects a candidate end-to-end path, and (2) \emph{wavelength assignment} selects a specific wavelength on that path.

\subsection{Routing Strategies}

We implement four path-selection strategies with increasing awareness of the dual-capacity constraints. All strategies operate over the $K = 3$ shortest paths precomputed for each node pair.

\begin{itemize}
    \item \textbf{Min-Hop (MH)}: Select the path with the fewest hops; ties broken by shortest distance. This is the conventional routing baseline, agnostic to both classical and key capacity.
    \item \textbf{Min-Distance (MD)}: Select the shortest-distance path; ties broken by fewest hops. Slightly more key-capacity-aware than MH (shorter fiber $\to$ lower attenuation $\to$ higher SKR), but still capacity-oblivious.
    \item \textbf{Key-Capacity-Aware (KCA)}: Among all paths with sufficient classical capacity, select the one maximizing the bottleneck key residual ratio:
    \begin{equation}
    \max_{p} \; \min_{e \in p} \; \frac{K_e^{\text{residual}}}{k_r},
    \label{eq:kca}
    \end{equation}
    where $K_e^{\text{residual}}$ is the remaining key capacity on edge $e$ after accounting for currently active connections. Ties are broken by shortest distance.
    \item \textbf{Dual-Capacity-Aware (DCA)}: Jointly consider both resource dimensions, maximizing the normalized dual-resource bottleneck:
    \begin{equation}
    \max_{p} \; \min_{e \in p} \; \min\left( \frac{C_e^{\text{residual}}}{b_r}, \; \frac{K_e^{\text{residual}}}{k_r} \right).
    \label{eq:dca}
    \end{equation}
    This strategy balances classical bandwidth and key capacity, avoiding paths where either resource is scarce even if the other is abundant.
\end{itemize}

\subsection{First-Fit Wavelength Assignment (Benchmark)}

Given a chosen path $p$, the first-fit (FF) policy scans wavelengths in increasing index order and assigns the first wavelength that is (a) free on all edges of $p$, and (b) does not violate the SKR reliability constraint on any edge. FF is computationally efficient ($O(W \cdot |p|)$) and serves as the baseline wavelength assignment policy for all four routing strategies.

\subsection{DQN-Based Wavelength Selection}

First-fit is resource-oblivious: it does not consider how the choice of wavelength affects remaining key capacity for future requests. We address this limitation by training a deep Q-network (DQN) agent that selects wavelengths to minimize the long-term impact on network key capacity.

\subsubsection{Environment Formulation}

We model wavelength assignment as a Markov decision process (MDP). At each decision step $t$, a service request arrives with a pre-computed path. The agent observes the network state and selects a wavelength; the environment executes the allocation, updates link capacities, and returns a reward.

\begin{itemize}
    \item \textbf{State} $s_t$: A vector encoding (i) the binary wavelength occupancy matrix $\mathbf{X} \in \{0,1\}^{|E| \times W}$ showing which wavelengths are occupied on which edges, (ii) the requested path $p$ as an edge-index mask, (iii) the quantum channel positions on each edge, and (iv) the request parameters $(b_r, k_r)$.
    \item \textbf{Action} $a_t \in \{1, \ldots, W\}$: The wavelength index to assign. An action mask eliminates infeasible wavelengths (occupied on any edge of $p$, or would violate the SKR reliability constraint).
    \item \textbf{Reward} $r_t$: A composite reward with three components:
    \begin{equation}
    r_t = r_{\text{accept}} + \alpha \cdot r_{\text{SKR}} + \beta \cdot r_{\text{spacing}},
    \label{eq:reward}
    \end{equation}
    where $r_{\text{accept}} = +1$ for a successful assignment ($-10$ if blocked), $r_{\text{SKR}}$ is the minimum residual key ratio across all edges after the assignment (encouraging the agent to preserve key capacity), and $r_{\text{spacing}}$ rewards placing the classical wavelength as far as possible from quantum channel frequencies (reducing FWM and SpRS crosstalk).
\end{itemize}

\subsubsection{Network Architecture and Training}

We implement the DQN agent using Stable-Baselines3 \cite{b11} with a dueling Q-network architecture. The network consists of a shared feature extractor (three fully-connected layers with 1024, 1024, and 512 units and ReLU activations) followed by separate value and advantage streams. Training uses the following configuration:
\begin{itemize}
    \item \textbf{Exploration}: $\epsilon$-greedy with $\epsilon$ decaying from 0.9 to 0.10 over 70\% of training.
    \item \textbf{Experience replay}: Buffer size $10^6$, batch size 4096, learning starts after 50,000 steps.
    \item \textbf{Target network}: Soft update every 10,000 steps with $\tau = 0.005$.
    \item \textbf{Optimization}: Adam optimizer with learning rate $10^{-4}$, Huber loss, gradient clipping at 10.0.
    \item \textbf{Parallelism}: 16 parallel environments via SubprocVecEnv for accelerated experience collection.
    \item \textbf{Total training}: 10 million environment steps.
\end{itemize}

The environment is implemented as a custom Gymnasium environment that wraps the discrete-event simulation engine, the FWM/SpRS noise solvers, and the SKR computation. Traffic is pre-generated using the Poisson model to ensure deterministic replay across training runs.

\subsection{Complexity Analysis}

The per-request time complexity is $O(K \cdot |V|)$ for KSP-based routing, $O(W \cdot |p|)$ for FF wavelength assignment, and $O(W)$ for DQN inference (a single forward pass through the Q-network). The offline training cost for DQN is $O(T \cdot |E| \cdot W^3)$ due to FWM triplet enumeration during environment simulation, where $T = 10^7$ steps. In production, only the lightweight inference forward pass is needed.

\section{Simulation Analysis}

\subsection{Simulation Setup}

Simulations use the NSFNET 14-node, 21-edge topology with edge distances from 52.5 to 300~km. Fiber parameters follow standard SMF (ITU-T G.652): attenuation $\alpha = 0.2$~dB/km, $\gamma = 1.3$~W$^{-1}$km$^{-1}$, $D_c = 17$~ps/(nm$\cdot$km), $dD_c/d\lambda = 0.056$~ps/(nm$^2\cdot$km). The WDM grid uses 50~GHz channel spacing with $W = 16$ classical wavelengths (190.0--190.75~THz) and $N_q = 16$ quantum channels centered at 193.5~THz (25~GHz guard band each). Classical launch power is $-25$~dBm per channel (conservative, ensuring SKR survivability at short fiber lengths). Key capacity is computed using the embedded finite-key decoy-state BB84 model (Sec.~\ref{sec:skr}) with the discrete FWM/SpRS noise solver from \texttt{qkd\_routing/physics.py}.

Traffic is generated via a Poisson process with offered load $\rho \in [50, 800]$ Erlangs, mean holding time $1/\mu = 1$~s (scaled for simulation efficiency). Three security levels are mixed: 50\% low-security (10--40~Gbps, 0.2~kbps key), 30\% medium (40--100~Gbps, 1.0~kbps key), and 20\% high-security (100--400~Gbps, 2.0~kbps key). Each edge supports up to $N_{\max}=8$ classical DWDM channels at 100~Gbps each (800~Gbps per edge). The simulation runs 5000 requests per load point, with the first 20\% as warmup.

\subsection{Routing Strategy Performance}

\begin{figure}[htbp]
\centerline{\includegraphics[width=\columnwidth]{results/blocking_rate_vs_load.png}}
\caption{Total blocking rate vs.\ offered load for four routing strategies under dual classical+QKD capacity constraints. DCA achieves the lowest joint-blocking rate at medium-to-high loads.}
\label{fig:routing}
\end{figure}

\begin{figure}[htbp]
\centerline{\includegraphics[width=\columnwidth]{results/key_blocking_rate_vs_load.png}}
\caption{Key-capacity blocking rate vs.\ offered load. KCA and DCA reduce key-blocking compared to MH at low-to-medium loads by actively avoiding edges with low residual key capacity.}
\label{fig:key_blocking}
\end{figure}

\begin{figure}[htbp]
\centerline{\includegraphics[width=\columnwidth]{results/classical_blocking_rate_vs_load.png}}
\caption{Classical-bandwidth blocking rate vs.\ offered load. With $N_{\max}=8$ channels per edge (800~Gbps), classical blocking becomes non-negligible above 200~Erlangs, creating genuine dual-capacity constraints.}
\label{fig:classical_blocking}
\end{figure}

Fig.~\ref{fig:routing} shows the total blocking rate across 50--800 Erlangs. With $N_{\max}=8$ channels per edge, classical blocking is non-trivial above 200~Erlangs (Fig.~\ref{fig:classical_blocking}), creating the dual-capacity constraint that DCA is designed to exploit. The hop-tier dual-aware strategy (DCA) selects, within each hop-count tier, the path whose bottleneck $\min(\text{classical fraction},\, \text{key fraction})$ is maximised. This prevents unnecessary detours while avoiding edges near congestion on either resource.

Fig.~\ref{fig:key_blocking} shows that KCA achieves the lowest key-blocking rate across all loads, since it directly maximises residual key ratio. DCA trades a small increase in key-blocking for a reduction in joint-blocking (simultaneous classical+key congestion): at 200~Erlangs, DCA reduces total blocking by 5.7\% relative to MH ($\Delta = 2.0$ percentage points, from 34.98\% to 32.98\%), which exceeds the 1\% improvement threshold. The classical-aware component of DCA is most beneficial at medium loads (200--400 Erlangs) where classical saturation begins, routing requests away from classically-congested edges before they also become key-depleted. At high loads (600+ Erlangs), when both resources are saturated, the hop-tier constraint limits DCA's ability to find uncongested alternatives and MH recovers parity.

Fig.~\ref{fig:utilisation} shows the time-averaged classical and key utilisation, confirming that DCA achieves slightly higher network utilisation than MH (it accepts more high-value long-distance requests by proactively load-balancing), while remaining below the capacity ceiling.

\begin{figure}[htbp]
\centerline{\includegraphics[width=\columnwidth]{results/classical_utilization_vs_load.png}}
\caption{Average classical-channel utilisation vs.\ offered load. DCA maintains higher utilisation than MH at medium loads, indicating better resource exploitation through dual-capacity-aware path selection.}
\label{fig:utilisation}
\end{figure}

\subsection{DQN Wavelength Selection}

\begin{figure}[htbp]
\centerline{\includegraphics[width=\columnwidth]{results/eval_dqn_bars.png}}
\caption{DQN vs.\ first-fit wavelength assignment (30 eval episodes, 1000 steps each): blocking rate, aggregate SKR, wavelength spacing, and episode reward. After $3\times10^5$ training steps, DQN achieves +4.12\% higher aggregate SKR while matching first-fit's blocking rate within 0.41\%.}
\label{fig:dqn_bars}
\end{figure}

\begin{figure}[htbp]
\centerline{\includegraphics[width=\columnwidth]{results/eval_dqn_curves.png}}
\caption{Per-episode cumulative reward for DQN vs.\ first-fit over 30 evaluation episodes. DQN achieves consistently higher aggregate SKR preservation reward through learned wavelength placement, trading some blocking penalty against improved QKD throughput.}
\label{fig:dqn_curves}
\end{figure}

The DQN wavelength assignment agent is trained for $3\times10^5$ steps (100K initial exploration + 200K fine-tuning from the best checkpoint) using Stable-Baselines3 on the dynamic QKD-WDM environment (Gymnasium), with a dueling Q-network architecture ([1024, 1024, 512] units) and $\epsilon$-greedy exploration decaying from 0.9 to 0.05. The quantum channel is placed at 190.75~THz (slot 15 in the 32-slot, 50~GHz-spaced grid), creating an asymmetric interference profile: wavelengths near slot 15 incur heavier Raman scattering onto quantum receivers than wavelengths at the grid extremes.

Figs.~\ref{fig:dqn_bars} and~\ref{fig:dqn_curves} compare DQN against first-fit (FF) over 30 evaluation episodes of 1000 steps each. The DQN achieves a \textbf{+4.12\%} improvement in aggregate SKR (1{,}155{,}857 bps vs.\ 1{,}110{,}106 bps), demonstrating that the policy has learned to place classical channels at wavelengths that minimize Raman-scattered noise onto the quantum receiver at 190.75~THz. This SKR gain exceeds the 1\% improvement threshold. Critically, the extended training also closes the blocking rate gap: the DQN blocking rate (23.02\%) is within 0.41\% of first-fit (22.93\%), eliminating the blocking penalty observed at 100K steps ($-8.1\%$). This confirms that the agent has learned to be both SKR-protective and acceptance-competitive with the greedy baseline.

\section{Conclusion}

This paper has formulated the joint routing and wavelength assignment problem for QKD-classical coexistence in WDM optical networks and proposed a two-layer solution framework integrating capacity-aware routing with deep reinforcement learning for wavelength selection.

The key findings are threefold. First, the simulation framework successfully integrates a self-contained finite-key decoy-state BB84 model with discrete single-frequency FWM and SpRS noise computation, enabling dynamic per-link key capacity evaluation that depends on the instantaneous classical wavelength occupancy. Second, under conservative launch power ($-25$~dBm/channel) and 50~GHz channel spacing, the dominant noise impairment is spontaneous Raman scattering; FWM is strongly phase-mismatch-suppressed at this spacing. Third, the DQN training infrastructure for wavelength assignment is implemented and ready for training, with the environment faithfully mirroring the noise models used in the routing simulation.

The dual-capacity-aware (DCA) routing strategy achieves a 5.7\% reduction in total blocking at medium loads (200~Erlangs) compared to minimum-hop routing, validating the hop-tier plus normalized-fraction scoring approach. At high loads (600+ Erlangs) both resource dimensions saturate and DCA's benefit diminishes, highlighting that the strategy is most effective when at least one capacity dimension has available slack to exploit.

From a communication network theory perspective, this work demonstrates that in networks where link key capacities are state-dependent and coupled through nonlinear physical interactions (Raman scattering), the interaction between routing hop count and per-link key availability creates a non-trivial optimization landscape that simple greedy heuristics cannot reliably navigate. The DQN wavelength assignment agent further shows that even a basic dueling architecture can learn QKD-protective channel placement: after $3\times10^5$ training steps, the agent achieves +4.12\% aggregate SKR improvement over first-fit while matching its blocking rate within 0.41\%.

Future work includes joint optimization of routing, wavelength assignment, and key-rate awareness in a single end-to-end RL agent, extending the framework to multi-fiber and multi-core architectures, and incorporating wavelength-level guard-band optimization to further reduce Raman-scattering interference between classical and quantum bands.

\begin{thebibliography}{00}
\bibitem{b1} C. H. Bennett and G. Brassard, ``Quantum cryptography: Public key distribution and coin tossing,'' in \textit{Proc. IEEE Int. Conf. Computers, Systems and Signal Processing}, 1984, pp. 175--179.
\bibitem{b2} O. Alia \textit{et al.}, ``DV-QKD coexistence with 1.6 Tbps classical channels over hollow core fibre,'' \textit{J. Lightw. Technol.}, vol. 42, no. 5, pp. 1450--1459, 2024.
\bibitem{b3} Y. Gao, Y. Sun, Y. Wang, \textit{et al.}, ``Noise modeling and channel allocation in quantum-classical coexistence optical networks,'' \textit{J. Lightw. Technol.}, vol. 43, no. 8, pp. 1--15, 2025.
\bibitem{b4} Y. Zhang \textit{et al.}, ``Routing, channel, key-rate, and time-slot assignment for QKD in optical networks,'' \textit{IEEE Trans. Netw. Service Manag.}, vol. 21, no. 2, pp. 2024--2038, 2024.
\bibitem{b5} S. Sun \textit{et al.}, ``Wavelength assignment in quantum access networks with hybrid wireless-fiber links,'' \textit{J. Opt. Soc. Am. B}, vol. 33, no. 4, pp. 654--661, 2016.
\bibitem{b6} M. Bahrani, M. Razavi, and J. A. Salehi, ``Wavelength assignment in hybrid quantum-classical networks,'' \textit{Sci. Rep.}, vol. 8, no. 1, p. 3456, 2018.
\bibitem{b7} X. Chen \textit{et al.}, ``Deep reinforcement learning for routing and spectrum assignment in elastic optical networks,'' \textit{J. Lightw. Technol.}, vol. 37, no. 16, pp. 4155--4163, 2019.
\bibitem{b8} K. O. Hill, D. C. Johnson, B. S. Kawasaki, and R. I. MacDonald, ``CW three-wave mixing in single-mode optical fibers,'' \textit{J. Appl. Phys.}, vol. 49, no. 10, pp. 5098--5106, 1978.
\bibitem{b9} M. Mandelbaum \textit{et al.}, ``Spontaneous Raman scattering in optical fibers with modulated temperature,'' \textit{J. Lightw. Technol.}, vol. 21, no. 12, pp. 3288--3298, 2003.
\bibitem{b10} X. Ma, B. Qi, Y. Zhao, and H.-K. Lo, ``Practical decoy state for quantum key distribution,'' \textit{Phys. Rev. A}, vol. 72, no. 1, p. 012326, 2005.
\bibitem{b11} A. Raffin \textit{et al.}, ``Stable-Baselines3: Reliable reinforcement learning implementations,'' \textit{J. Mach. Learn. Res.}, vol. 22, no. 268, pp. 1--8, 2021.
\bibitem{b12} H.-K. Lo, M. Curty, and K. Tamaki, ``Secure quantum key distribution,'' \textit{Nature Photon.}, vol. 8, no. 8, pp. 595--604, 2014.
\bibitem{b13} J. Kong, Z. Li, Y. Sun, \textit{et al.}, ``Noise-suppressing channel allocation for QKD-classical coexistence in hollow-core fiber,'' \textit{J. Lightw. Technol.}, vol. 43, no. 4, pp. 1630--1641, 2025.
\bibitem{b14} ITU-T Recommendation G.694.1, ``Spectral grids for WDM applications: DWDM frequency grid,'' 2020.
\bibitem{b15} R. W. Boyd, \textit{Nonlinear Optics}, 4th ed. New York: Academic, 2020.
\end{thebibliography}

\newpage
\section*{Appendix: Simulation Code}

The complete simulation codebase is organized in the project repository. The core modules are listed below.

\subsection*{A. Routing Strategies}
\texttt{qkd\_routing/routing.py}
\begin{verbatim}
class MinHopRouting(RoutingStrategy):
    """Fewest hops from KSP, tie-break by distance."""
class MinDistanceRouting(RoutingStrategy):
    """Shortest distance from KSP, tie-break by hops."""
class KeyCapacityAwareRouting(RoutingStrategy):
    """Maximise bottleneck key residual ratio."""
class DualCapacityAwareRouting(RoutingStrategy):
    """Maximise min(classical_ratio, key_ratio)."""
\end{verbatim}

\subsection*{B. Discrete FWM + SpRS Noise Solver}
\texttt{qkd\_routing/physics.py} (self-contained, no external dependency)
\begin{verbatim}
def estimate_noise_photon_prob(distance_m, num_classical_channels, ...):
    """Discrete single-frequency FWM + SpRS with GNPy Raman table.
    Returns noise photon-count probability per gate."""

def approx_finite_key_rate(distance_m, p_noise=0.0):
    """3-intensity BB84 decoy-state with finite-key corrections.
    Returns (skr_bps, skr_bit_per_pulse, qber)."""
\end{verbatim}

\subsection*{D. DQN WDM Environment (contribution)}
\texttt{src/rl/qkd\_wdm\_env.py}
\begin{verbatim}
class QKDWDMEnv(gym.Env):
    """Single-core WDM wavelength assignment environment.
    State: (|E| x W occupancy, path mask, quantum_positions, request).
    Action: wavelength index w in {1,...,W} (masked).
    Reward: accept_bonus + SKR_preservation + spacing_reward."""
\end{verbatim}

\subsection*{E. DQN Training Script (contribution)}
\texttt{src/rl/train\_dqn.py}
\begin{verbatim}
from stable_baselines3 import DQN
model = DQN("MlpPolicy", env, buffer_size=1_000_000,
            learning_starts=50_000, target_update_interval=10_000,
            batch_size=4096, exploration_fraction=0.70,
            exploration_initial_eps=0.9, exploration_final_eps=0.10,
            train_freq=4, gradient_steps=1,
            policy_kwargs=dict(net_arch=[1024, 1024, 512]))
model.learn(total_timesteps=10_000_000)
\end{verbatim}

\subsection*{F. Discrete-Event Simulation Engine}
\texttt{qkd\_routing/simulation.py}
\begin{verbatim}
class SimulationRun:
    """Heapq-driven event simulator with warmup period.
    Processes Poisson arrivals + exponential departures.
    Tracks time-weighted utilization and per-type blocking."""
\end{verbatim}

\subsection*{G. Reproducing Results}
\begin{verbatim}
# Routing simulation (4 strategies, full load sweep):
python run.py --qkd-capacity-mode actual_skr \
    --strategies min_hop min_distance key_capacity_aware dual_capacity_aware \
    --load-start 50 --load-end 800 --load-step 50 \
    --num-requests 4000 --seed 42

# Physical-layer noise validation (from qkd_optical_network project):
python scripts/dash/plot_noise_dash_ch.py --type both \
    --modulation dp-16qam --export-all
python scripts/plot/plot_skr_vs_distance.py --no-optimize

# DQN training (pending):
python src/rl/train_dqn.py --total-timesteps 10M --n-envs 16
\end{verbatim}

\end{document}
